\documentclass[12pt,a4paper]{article}

\usepackage{fontspec}
\usepackage{polyglossia}
\setmainlanguage{vietnamese}
\setmainfont{DejaVu Serif}
\setsansfont{DejaVu Sans}
\setmonofont[Scale=0.86]{DejaVu Sans Mono}

\usepackage{amsmath,amssymb,amsthm,mathtools,bm}
\usepackage{geometry}
\usepackage{graphicx}
\usepackage{booktabs,array,multirow,tabularx,longtable}
\usepackage{xcolor}
\usepackage[colorlinks=true,linkcolor=blue!65!black,citecolor=green!45!black,urlcolor=blue!70!black]{hyperref}
\usepackage{enumitem}
\usepackage{tikz}
\usetikzlibrary{positioning,arrows.meta,shapes.geometric,calc,fit,backgrounds}
\usepackage{pgfplots}
\pgfplotsset{compat=1.18}
\usepackage{caption}
\usepackage{subcaption}
\usepackage{float}
\usepackage{fancyhdr}
\usepackage{titlesec}
\usepackage{tcolorbox}
\usepackage{microtype}

\geometry{top=2.45cm,bottom=2.45cm,left=3cm,right=2.45cm}
\setlist[itemize]{topsep=3pt,itemsep=2pt,leftmargin=1.4em}
\setlist[enumerate]{topsep=3pt,itemsep=2pt,leftmargin=1.5em}
\renewcommand{\arraystretch}{1.18}

\definecolor{mainblue}{HTML}{0A3D62}
\definecolor{softblue}{HTML}{EAF2F8}
\definecolor{deepgreen}{HTML}{1E8449}
\definecolor{warnorange}{HTML}{B9770E}
\definecolor{darkred}{HTML}{922B21}
\definecolor{softgray}{HTML}{F4F6F7}

\pagestyle{fancy}
\fancyhf{}
\lhead{\textcolor{gray!75}{\small Báo cáo RELGT}}
\rhead{\textcolor{gray!75}{\small RELGT -- RELGT++ -- RELGT-X}}
\cfoot{\thepage}
\renewcommand{\headrulewidth}{0.35pt}

\titleformat{\section}{\Large\bfseries\color{mainblue}}{\thesection}{0.75em}{}[\color{mainblue!35}\titlerule]
\titleformat{\subsection}{\large\bfseries\color{mainblue!85!black}}{\thesubsection}{0.65em}{}
\titleformat{\subsubsection}{\normalsize\bfseries\color{black!80}}{\thesubsubsection}{0.55em}{}

\newtheorem{definition}{Định nghĩa}[section]
\newtheorem{proposition}{Mệnh đề}[section]
\newtheorem{remark}{Nhận xét}[section]

\tcbuselibrary{skins,breakable}
\newtcolorbox{summarybox}[1][]{
  enhanced,breakable,
  colback=softblue,
  colframe=mainblue!65!black,
  fonttitle=\bfseries,
  title=#1
}
\newtcolorbox{insightbox}[1][]{
  enhanced,breakable,
  colback=green!4!white,
  colframe=deepgreen!70!black,
  fonttitle=\bfseries,
  title=#1
}
\newtcolorbox{warnbox}[1][]{
  enhanced,breakable,
  colback=orange!5!white,
  colframe=warnorange!80!black,
  fonttitle=\bfseries,
  title=#1
}
\newtcolorbox{mathbox}[1][]{
  enhanced,breakable,
  colback=gray!4!white,
  colframe=gray!55!black,
  fonttitle=\bfseries,
  title=#1
}

\newcommand{\RelGT}{\textsc{RelGT}}
\newcommand{\RelGTpp}{\textsc{RelGT++}}
\newcommand{\RelGTx}{\textsc{RelGT-X}}
\newcommand{\RDL}{\textsc{RDL}}
\newcommand{\REG}{\textsc{REG}}
\newcommand{\concat}{\mathbin{\Vert}}

\begin{document}

% ============================================================
% Trang bìa
% ============================================================
\begin{titlepage}
\centering
\vspace*{0.7cm}

{\small\scshape\color{gray!75} Báo cáo tổng hợp và phân tích kỹ thuật}\\[0.8em]
\rule{\textwidth}{2.4pt}\\[0.7em]
{\Huge\bfseries\color{mainblue} RELGT}\\[0.35em]
{\LARGE\bfseries Relational Graph Transformer}\\[0.45em]
{\Large\color{black!70} Từ \RelGT{} gốc đến \RelGTpp{} và đề xuất \RelGTx}\\[0.7em]
\rule{\textwidth}{1pt}\\[1.2em]

\begin{tcolorbox}[colback=softgray,colframe=gray!45,width=0.93\textwidth]
\centering
\textbf{Phạm vi báo cáo}\\[0.3em]
Phân tích chi tiết paper \RelGT{}, giải thích kiến trúc \RelGT{} gốc,
tổng hợp các cải tiến trong \RelGTpp{}, và trình bày định hướng mở rộng
\RelGTx{} cho dữ liệu quan hệ đa bảng có thời gian.
\end{tcolorbox}

\vspace{1.3em}
\begin{tabular}{ll}
\toprule
\textbf{Chủ đề} & Graph Transformer cho Relational Deep Learning \\
\textbf{Nguồn chính} & \texttt{relgt\_report\_vi.tex}, \texttt{relgtpp\_plusplus.tex}, \texttt{relgtx\_x.tex} \\
\textbf{Bài báo gốc} & Dwivedi et al., \textit{Relational Graph Transformer}, arXiv:2505.10960 \\
\textbf{Benchmark} & RelBench, 21 tác vụ trên cơ sở dữ liệu quan hệ \\
\textbf{Ngày lập báo cáo} & \today \\
\bottomrule
\end{tabular}

\vfill
\begin{tikzpicture}[
  box/.style={draw,rounded corners=5pt,minimum width=3.2cm,minimum height=1.05cm,
              align=center,font=\small\bfseries},
  arrow/.style={-Stealth,thick,mainblue}
]
\node[box,fill=blue!8,draw=mainblue] (a) {\RelGT\\Token hóa 5 phần};
\node[box,fill=green!8,draw=deepgreen,right=1.3cm of a] (b) {\RelGTpp\\Gated Fusion + Codebook};
\node[box,fill=orange!10,draw=warnorange,right=1.3cm of b] (c) {\RelGTx\\Schema + Causal + Meta};
\draw[arrow] (a) -- (b);
\draw[arrow] (b) -- (c);
\end{tikzpicture}

\vspace*{0.5cm}
\end{titlepage}

% ============================================================
% Trang mục lục
% ============================================================
\renewcommand{\contentsname}{Mục lục}
\clearpage
\thispagestyle{plain}
\tableofcontents
\clearpage

% ============================================================
\section{Giới thiệu}
% ============================================================

\begin{summarybox}[Ý tưởng và paper gốc]
\RelGT{} là kiến trúc Graph Transformer được thiết kế cho dữ liệu quan hệ
đa bảng, được giới thiệu trong paper gốc \textit{Relational Graph Transformer}
của Vijay Prakash Dwivedi, Sri Jaladi, Yangyi Shen, Federico López,
Charilaos I. Kanatsoulis, Rishi Puri, Matthias Fey và Jure Leskovec
(arXiv:2505.10960v1, 16/05/2025)~\cite{relgt2025}. Paper này đặt vấn đề rằng
dữ liệu doanh nghiệp thực tế thường không tồn tại như một bảng đơn lẻ, mà là
một hệ nhiều bảng, nhiều kiểu thực thể, nhiều quan hệ khóa chính--khóa ngoại
và nhiều mốc thời gian. Vì vậy, mô hình cần học đồng thời đặc trưng bảng,
cấu trúc quan hệ, ngữ nghĩa schema và động lực thời gian.

Khác với GNN truyền thống, \RelGT{} biến mỗi node trong đồ thị quan hệ thành
một token giàu ngữ cảnh gồm đặc trưng bảng, loại node, khoảng cách hop,
thời gian tương đối và mã hóa vị trí cục bộ. Sau đó mô hình dùng attention
cục bộ để học tương tác trong đồ thị con và attention toàn cục đến một tập
centroid để thu nhận ngữ cảnh xa.
\end{summarybox}

Báo cáo này không chỉ tóm tắt paper gốc, mà còn mở rộng phân tích theo ba
thế hệ kiến trúc:
\begin{itemize}
  \item \textbf{\RelGT{} gốc}: đóng góp nền tảng về token hóa 5 thành phần,
        lấy mẫu nhận biết thời gian, local transformer và global centroid.
  \item \textbf{\RelGTpp{}}: phiên bản triển khai/cải tiến, tăng tính ổn định
        và biểu đạt bằng CrossModalGatedFusion, SwiGLU, DropPath,
        Multi-View Readout, GumbelSoftCodebook và CrossAttentionBridge.
  \item \textbf{\RelGTx{}}: hướng nghiên cứu kế tiếp, đưa schema prior,
        temporal causal disentanglement, hierarchical codebook, relational
        positional encoding, meta-learning tokenization và contrastive schema
        alignment vào cùng một khung.
\end{itemize}

Nhìn tổng thể, \RelGT{} giải quyết một nút thắt rất thực tế: dữ liệu doanh
nghiệp không nằm trong một ma trận duy nhất mà trong nhiều bảng, nhiều kiểu
thực thể, nhiều quan hệ khóa chính--khóa ngoại và nhiều mốc thời gian. Bài
toán không chỉ là học đặc trưng tabular, cũng không chỉ là học graph; nó là
học trên một \emph{hệ quan hệ có cấu trúc, có thời gian và có ngữ nghĩa schema}.

% ============================================================
\section{Bối cảnh: Relational Deep Learning}
% ============================================================

\subsection{Dữ liệu quan hệ và giới hạn của cách tiếp cận cũ}

Trong các hệ thống thương mại điện tử, tài chính, y tế, mạng xã hội hoặc
vận hành doanh nghiệp, dữ liệu thường được lưu trong cơ sở dữ liệu quan hệ:
bảng người dùng, bảng sản phẩm, bảng đơn hàng, bảng đánh giá, bảng sự kiện,
bảng thanh toán, v.v. Các bảng này được nối bằng khóa chính và khóa ngoại.

Cách làm truyền thống thường gồm ba bước: chọn bảng mục tiêu, join nhiều bảng
liên quan, sinh đặc trưng thủ công, rồi huấn luyện một mô hình tabular như
LightGBM hoặc mạng neural cho bảng phẳng. Cách làm đó có ba hạn chế lớn:
\begin{enumerate}
  \item \textbf{Tốn công thiết kế đặc trưng}: cần chuyên gia hiểu nghiệp vụ,
        schema và cửa sổ thời gian.
  \item \textbf{Mất cấu trúc quan hệ}: sau khi join/aggregate, nhiều đường đi
        quan hệ trong schema bị nén thành thống kê đơn giản.
  \item \textbf{Dễ rò rỉ thời gian}: nếu aggregation dùng dữ liệu sau thời điểm
        dự báo, mô hình nhìn thấy tương lai.
\end{enumerate}

\RDL{} đề xuất học trực tiếp từ cơ sở dữ liệu quan hệ bằng cách chuyển nó
thành đồ thị không đồng nhất có thời gian. Khi đó mỗi hàng trong bảng là một
node, mỗi quan hệ khóa chính--khóa ngoại là edge, và mỗi bảng tương ứng với
một loại node.

\subsection{Relational Entity Graph}

\begin{definition}[Relational Entity Graph]
Một \REG{} được định nghĩa là:
\[
  \mathcal{G}=(\mathcal{V},\mathcal{E},\varphi,\psi,\tau),
\]
trong đó $\mathcal{V}$ là tập node/thực thể, $\mathcal{E}$ là tập cạnh quan hệ,
$\varphi:\mathcal{V}\rightarrow\mathcal{T}$ ánh xạ node về loại bảng,
$\psi:\mathcal{E}\rightarrow\mathcal{R}$ ánh xạ cạnh về loại quan hệ, và
$\tau:\mathcal{V}\rightarrow\mathbb{R}$ là thời điểm của node hoặc sự kiện.
\end{definition}

\REG{} có ba đặc điểm quan trọng:
\begin{itemize}
  \item \textbf{Schema-defined}: edge không tùy ý mà do schema cơ sở dữ liệu
        quy định.
  \item \textbf{Temporal}: dự báo phải tôn trọng thứ tự thời gian, đặc biệt
        trong các tác vụ như churn, fraud, sales forecasting.
  \item \textbf{Heterogeneous}: mỗi bảng có miền đặc trưng, kiểu dữ liệu và
        ý nghĩa nghiệp vụ riêng.
\end{itemize}

\subsection{Vì sao GNN truyền thống chưa đủ}

GNN học bằng message passing: sau $L$ lớp, một node nhận thông tin từ vùng
$L$-hop. Trên dữ liệu quan hệ, điều này có thể thiếu linh hoạt:
\begin{itemize}
  \item Phụ thuộc xa cần đi qua nhiều bảng trung gian nên dễ bị \emph{over-squashing}.
  \item Hai thực thể cùng kiểu có thể liên quan qua schema nhưng không có cạnh
        trực tiếp trong đồ thị con nhỏ.
  \item Heterogeneity và temporal dynamics thường bị nhét vào cùng một vector,
        làm mất ranh giới ngữ nghĩa giữa nội dung, kiểu node, thời gian và vị trí.
\end{itemize}

\begin{insightbox}[Trực giác]
GNN truyền thống giỏi lan truyền tín hiệu theo cạnh. \RelGT{} đặt câu hỏi khác:
nếu đã có một đồ thị con quanh seed node, ta có thể xem các node trong đó như
một tập token và để Transformer học mọi tương tác quan trọng giữa chúng không?
\end{insightbox}

% ============================================================
\section{\RelGT{} Gốc}
% ============================================================

\subsection{Tổng quan kiến trúc}

\RelGT{} gồm ba bước chính:
\[
  \text{REG}
  \xrightarrow{\text{lấy mẫu nhận biết thời gian}}
  \text{đồ thị con quanh seed}
  \xrightarrow{\text{token hóa 5 thành phần}}
  \text{local + global Transformer}
  \xrightarrow{}
  \hat{y}.
\]

\begin{figure}[H]
\centering
\begin{tikzpicture}[
  node distance=0.9cm and 1.0cm,
  box/.style={draw,rounded corners=5pt,minimum width=2.7cm,minimum height=0.8cm,
              align=center,font=\small},
  enc/.style={box,fill=blue!9,draw=mainblue},
  mod/.style={box,fill=green!8,draw=deepgreen},
  out/.style={box,fill=orange!10,draw=warnorange},
  arr/.style={-Stealth,thick}
]
\node[enc] (db) {CSDL quan hệ\\đa bảng};
\node[enc,right=of db] (reg) {\REG\\hetero-temporal};
\node[enc,right=of reg] (sample) {Lấy mẫu\\2-hop};
\node[mod,right=of sample] (token) {Token hóa\\5 thành phần};
\node[mod,below right=0.6cm and 1.0cm of token] (local) {Local\\Transformer};
\node[mod,above right=0.6cm and 1.0cm of token] (global) {Global\\Centroid};
\node[out,right=2.3cm of token] (pred) {Dự đoán\\$\hat{y}$};
\draw[arr] (db) -- (reg);
\draw[arr] (reg) -- (sample);
\draw[arr] (sample) -- (token);
\draw[arr] (token) -- (local);
\draw[arr] (token) -- (global);
\draw[arr] (local) -- (pred);
\draw[arr] (global) -- (pred);
\end{tikzpicture}
\caption{Pipeline khái quát của \RelGT{}.}
\end{figure}

\subsection{Lấy mẫu lân cận nhận biết thời gian}

Với seed node $v_i$, \RelGT{} lấy $K-1$ node lân cận trong phạm vi 2-hop,
nhưng chỉ giữ những node không xảy ra sau seed:
\[
  \mathcal{N}(v_i)=
  \{v_j: d(v_i,v_j)\le 2,\; \tau(v_j)\le \tau(v_i)\}.
\]
Nếu số lân cận không đủ, mô hình dùng lấy mẫu có hoàn lại; nếu không có lân
cận hợp lệ, mô hình dùng fallback token từ toàn đồ thị.

\begin{mathbox}[Quy tắc lấy mẫu]
\[
  \mathcal{S}(v_i)=
  \begin{cases}
    \{v_i\}\cup \text{Sample}(\mathcal{N}(v_i),K-1),
      & |\mathcal{N}(v_i)|\ge K-1,\\
    \{v_i\}\cup \text{SampleWithReplacement}(\mathcal{N}(v_i),K-1),
      & 0<|\mathcal{N}(v_i)|<K-1,\\
    \{v_i\}\cup \text{FallbackSample}(\mathcal{V},K-1),
      & |\mathcal{N}(v_i)|=0.
  \end{cases}
\]
\end{mathbox}

Ràng buộc $\tau(v_j)\le\tau(v_i)$ là yếu tố rất quan trọng. Nó biến quá trình
lấy mẫu từ một thao tác graph thông thường thành một thao tác phù hợp cho
dự báo theo thời gian.

\subsection{Token hóa 5 thành phần}

Mỗi token $t(v_j\mid v_i)$ trong ngữ cảnh seed $v_i$ được biểu diễn bằng:
\[
  t(v_j\mid v_i)=
  \bigl(x_{v_j},\varphi(v_j),p(v_i,v_j),
  \tau(v_j)-\tau(v_i),\text{GNN-PE}_{v_j}\bigr).
\]

\begin{table}[H]
\centering
\caption{Năm thành phần token trong \RelGT{}.}
\begin{tabularx}{\textwidth}{clX}
\toprule
\textbf{\#} & \textbf{Thành phần} & \textbf{Vai trò} \\
\midrule
1 & Node feature $x_{v_j}$ & Nội dung bảng: số, phân loại, timestamp, văn bản/embedding. \\
2 & Node type $\varphi(v_j)$ & Cho biết token đến từ bảng nào, giúp mô hình phân biệt ngữ nghĩa thực thể. \\
3 & Hop distance $p(v_i,v_j)$ & Mô tả khoảng cách cấu trúc so với seed: 1-hop, 2-hop hoặc fallback. \\
4 & Relative time $\Delta\tau$ & Giữ thông tin thứ tự và khoảng cách thời gian giữa token và seed. \\
5 & Subgraph PE & Mã hóa vị trí cục bộ bằng GNN nhẹ trên đồ thị con. \\
\bottomrule
\end{tabularx}
\end{table}

Các thành phần được mã hóa riêng:
\begin{align}
  h_\text{feat}(v_j) &= \text{MultiModalEncoder}(x_{v_j}),\\
  h_\text{type}(v_j) &= W_\text{type}[\varphi(v_j)],\\
  h_\text{hop}(v_i,v_j) &= W_\text{hop}[p(v_i,v_j)],\\
  h_\text{time}(v_i,v_j) &= W_\text{time}\,\text{PE}(\tau(v_j)-\tau(v_i)),\\
  h_\text{pe}(v_j) &= \text{GNN}(A_\text{local},Z_\text{random})_j.
\end{align}

Sau đó:
\[
  h_\text{token}(v_j)=
  O\,[h_\text{feat}\concat h_\text{type}\concat h_\text{hop}
  \concat h_\text{time}\concat h_\text{pe}],
\]
với $O$ là ma trận trộn có thể học.

\subsection{Encoder đặc trưng bảng}

Đặc trưng bảng trong dữ liệu quan hệ thường đa phương thức: số, phân loại,
multi-categorical, timestamp, text embedding. \RelGT{} dùng encoder theo kiểu
dữ liệu, rồi đưa về cùng chiều $d$:
\[
  h_\text{feat}^{(t)}=\text{ResNet}^{(t)}
  \bigl(\text{Enc}_\text{num},\text{Enc}_\text{cat},
  \text{Enc}_\text{text},\text{Enc}_\text{time}\bigr).
\]
Điểm đáng chú ý là encoder được tách theo loại bảng $t$. Điều này phù hợp vì
cùng một cột số trong bảng sản phẩm và bảng người dùng không nhất thiết có
cùng phân phối hoặc cùng ý nghĩa.

\subsection{Time encoder}

\RelGT{} dùng positional encoding kiểu Transformer cho thời gian tương đối:
\[
  \text{PE}(\Delta\tau)_{2k}
  =\sin\left(\frac{\Delta\tau}{10000^{2k/d}}\right),\quad
  \text{PE}(\Delta\tau)_{2k+1}
  =\cos\left(\frac{\Delta\tau}{10000^{2k/d}}\right).
\]
Sau đó:
\[
  h_\text{time}=W_\text{time}\text{PE}(\Delta\tau)+b_\text{time}.
\]
Với token không có thời gian hợp lệ, triển khai có thể dùng một mask vector
có thể học để không ép giá trị thiếu thành một thời điểm giả.

\subsection{GNN positional encoding}

GNN-PE chạy một GNN nhẹ trên đồ thị con:
\[
  h_\text{pe}(v_j)=\text{GNN}(A_\text{local},Z_\text{random})_j.
\]
$Z_\text{random}$ được lấy ngẫu nhiên để phá đối xứng giữa các vị trí cấu trúc.
Về trực giác, nếu hai node có feature giống nhau nhưng nằm ở vị trí cấu trúc
khác nhau trong đồ thị con, GNN-PE tạo thêm tín hiệu giúp Transformer phân biệt.

\begin{remark}
GNN-PE là cầu nối giữa GNN và Transformer. Nó không dùng GNN làm bộ dự báo
chính, nhưng dùng GNN để cung cấp tọa độ cấu trúc cục bộ cho token.
\end{remark}

\subsection{Local Transformer}

Trên $K$ token của seed, \RelGT{} dùng self-attention:
\[
  \text{Attn}(Q,K,V)=
  \text{softmax}\left(\frac{QK^\top}{\sqrt{d_k}}\right)V.
\]
Vì $K$ cố định, local attention có chi phí $\mathcal{O}(K^2d)$ cho mỗi seed.
Sau các lớp Transformer, mô hình dùng readout để thu được biểu diễn cục bộ:
\[
  h_\text{local}(v_i)=
  \sum_{j=1}^{K}\alpha_j h_j,\quad
  \alpha_j=\frac{\exp(w^\top[h_i\concat h_j])}
  {\sum_k\exp(w^\top[h_i\concat h_k])}.
\]

\subsection{Global centroid module}

Local Transformer chỉ nhìn vào đồ thị con. Để có ngữ cảnh toàn cục, \RelGT{}
duy trì một codebook gồm $B$ centroid. Seed representation truy vấn vào
codebook:
\[
  h_\text{global}(v_i)=
  \text{Attn}(q(v_i),K_\text{cb},V_\text{cb}).
\]

Mỗi node được gán vào centroid gần nhất:
\[
  c_{v_i}=\arg\min_{b\in[B]}\|q_{v_i}-e_b\|_2.
\]
Codebook có thể cập nhật bằng EMA:
\begin{align}
  N_b^{(t)} &=
  \gamma N_b^{(t-1)}+(1-\gamma)\sum_i\mathbf{1}[c_{v_i}=b],\\
  e_b^{(t)} &=
  \frac{\gamma N_b^{(t-1)}e_b^{(t-1)}
  +(1-\gamma)\sum_{i:c_{v_i}=b}q_{v_i}}{N_b^{(t)}}.
\end{align}

\subsection{Kết hợp local và global}

Biểu diễn cuối cùng thường được lấy bằng:
\[
  h(v_i)=\text{MLP}\bigl([h_\text{local}(v_i)\concat h_\text{global}(v_i)]\bigr),
\quad
  \hat{y}_i=\text{Head}(h(v_i)).
\]
Tùy tác vụ, loss có thể là BCEWithLogits cho phân loại nhị phân/đa nhãn hoặc
MAE/MSE cho hồi quy.

\subsection{Độ phức tạp}

\begin{table}[H]
\centering
\caption{Độ phức tạp chính trong \RelGT{}.}
\begin{tabularx}{\textwidth}{lccX}
\toprule
\textbf{Thành phần} & \textbf{Time} & \textbf{Space} & \textbf{Ghi chú} \\
\midrule
Lấy mẫu & $\mathcal{O}(K)$ & $\mathcal{O}(K)$ & $K$ cố định theo cấu hình. \\
Token hóa & $\mathcal{O}(Kd)$ & $\mathcal{O}(Kd)$ & Gồm feature/type/hop/time/PE. \\
Local attention & $\mathcal{O}(K^2d)$ & $\mathcal{O}(K^2)$ & Chi phí cao nhất khi $K$ lớn. \\
Global attention & $\mathcal{O}(BHd_h)$ & $\mathcal{O}(B d)$ & $B$ là số centroid. \\
Head dự báo & $\mathcal{O}(d^2)$ & $\mathcal{O}(d)$ & Nhỏ so với attention/encoder. \\
\bottomrule
\end{tabularx}
\end{table}

\begin{insightbox}[Điểm cân bằng]
\RelGT{} chấp nhận chi phí $\mathcal{O}(K^2)$ trong local attention để đổi
lấy tương tác all-pair giữa các token trong đồ thị con. Vì $K$ cố định, chi
phí không tăng trực tiếp theo tổng số node toàn đồ thị.
\end{insightbox}

\subsection{Đánh giá trên RelBench}

RelBench gồm nhiều cơ sở dữ liệu và 21 tác vụ, bao gồm phân loại nhị phân,
hồi quy và phân loại đa nhãn. Các chỉ số thường dùng:
\begin{itemize}
  \item ROC-AUC cho phân loại nhị phân.
  \item MAE cho hồi quy.
  \item Macro AUPRC cho phân loại đa nhãn.
\end{itemize}

\begin{table}[H]
\centering
\small
\caption{Kết quả thực nghiệm của \RelGT{} trên các tác vụ RelBench tiêu biểu.
Bảng được tổng hợp từ paper gốc~\cite{relgt2025}. Với MAE, giá trị thấp hơn
là tốt hơn; với AUC, giá trị cao hơn là tốt hơn.}
\begin{subtable}[t]{0.49\textwidth}
\centering
\caption{MAE cho hồi quy thực thể.}
\resizebox{\linewidth}{!}{%
\begin{tabular}{llrrrrr}
\toprule
\textbf{Dataset} & \textbf{Task} & \textbf{RDL} & \textbf{HGT} &
\textbf{HGT+PE} & \textbf{\RelGT{}} & \textbf{Gain(\%)} \\
\midrule
rel-f1 & driver-position & 4.022 & 4.2263 & 4.3921 & \textbf{3.9170} & 2.61 \\
rel-avito & ad-ctr & 0.041 & 0.0462 & 0.0483 & \textbf{0.0345} & 15.85 \\
rel-event & user-attendance & 0.258 & 0.2635 & 0.2611 & \textbf{0.2502} & 2.79 \\
rel-trial & study-adverse & 44.473 & 45.1692 & \textbf{42.6484} & 43.9923 & 1.08 \\
 & site-success & 0.400 & 0.4428 & 0.4396 & \textbf{0.3263} & 18.43 \\
rel-amazon & user-ltv & 14.313 & 15.4120 & 15.8643 & \textbf{14.2665} & 0.32 \\
 & item-ltv & 50.053 & 55.8683 & 55.8493 & \textbf{48.9222} & 2.26 \\
rel-stack & post-votes & \textbf{0.065} & 0.0679 & 0.0680 & 0.0654 & \textcolor{darkred}{-0.62} \\
rel-hm & item-sales & 0.056 & 0.0641 & 0.0639 & \textbf{0.0536} & 4.29 \\
\bottomrule
\end{tabular}}
\end{subtable}
\hfill
\begin{subtable}[t]{0.49\textwidth}
\centering
\caption{AUC cho phân loại thực thể.}
\resizebox{\linewidth}{!}{%
\begin{tabular}{llrrrrr}
\toprule
\textbf{Dataset} & \textbf{Task} & \textbf{RDL} & \textbf{HGT} &
\textbf{HGT+PE} & \textbf{\RelGT{}} & \textbf{Gain(\%)} \\
\midrule
rel-f1 & driver-dnf & 0.7262 & 0.7077 & 0.7117 & \textbf{0.7587} & 4.48 \\
 & driver-top3 & 0.7554 & 0.7075 & 0.7627 & \textbf{0.8352} & 10.56 \\
rel-avito & user-clicks & 0.6590 & 0.6376 & 0.6457 & \textbf{0.6830} & 3.64 \\
 & user-visits & 0.6620 & 0.6432 & 0.6495 & \textbf{0.6678} & 0.88 \\
rel-event & user-repeat & \textbf{0.7689} & 0.6496 & 0.6536 & 0.7609 & \textcolor{darkred}{-1.04} \\
 & user-ignore & 0.8162 & \textbf{0.8247} & 0.8161 & 0.8157 & \textcolor{darkred}{-0.06} \\
rel-trial & study-outcome & 0.6860 & 0.5837 & 0.5921 & \textbf{0.6861} & 0.01 \\
rel-amazon & user-churn & \textbf{0.7042} & 0.6643 & 0.6619 & 0.7039 & \textcolor{darkred}{-0.04} \\
 & item-churn & \textbf{0.8281} & 0.7797 & 0.7803 & 0.8255 & \textcolor{darkred}{-0.31} \\
rel-stack & user-engagement & 0.9021 & 0.8847 & 0.8817 & \textbf{0.9053} & 0.35 \\
 & user-badge & \textbf{0.8986} & 0.8608 & 0.8566 & 0.8632 & \textcolor{darkred}{-3.94} \\
rel-hm & user-churn & \textbf{0.6988} & 0.6695 & 0.6569 & 0.6927 & \textcolor{darkred}{-0.87} \\
\bottomrule
\end{tabular}}
\end{subtable}
\end{table}

Kết quả trong paper cho thấy \RelGT{} có lợi thế so với các baseline GNN và
tabular trong nhiều tác vụ vì nó đồng thời khai thác feature bảng, schema,
thời gian, cấu trúc cục bộ và ngữ cảnh toàn cục.

% ============================================================
\section{\RelGTpp{}: Phiên bản Cải tiến}
% ============================================================

\subsection{Mục tiêu thiết kế}

\RelGTpp{} giữ pipeline của \RelGT{} nhưng thay nhiều khối bằng module mạnh
hơn. Nếu \RelGT{} trả lời câu hỏi ``làm sao đưa Transformer vào dữ liệu quan
hệ'', thì \RelGTpp{} trả lời câu hỏi thực dụng hơn: ``làm sao huấn luyện kiến
trúc đó ổn định, giàu biểu đạt và phù hợp với code production hơn''.

\begin{figure}[H]
\centering
\begin{tikzpicture}[
  box/.style={draw,rounded corners=4pt,minimum width=2.08cm,minimum height=0.58cm,
              align=center,font=\scriptsize},
  inp/.style={box,fill=gray!12,draw=gray!50},
  enc/.style={box,fill=blue!9,draw=mainblue},
  mod/.style={box,fill=green!8,draw=deepgreen},
  out/.style={box,fill=orange!10,draw=warnorange},
  arr/.style={-Stealth,thick}
]
\node[inp] (types) at (0,2.0) {node types};
\node[inp] (hops)  at (0,1.0) {hop ids};
\node[inp] (times) at (0,0.0) {rel times};
\node[inp] (feat)  at (0,-1.0) {table features};
\node[inp] (edge)  at (0,-2.0) {subgraph edges};

\node[enc] (et)  at (2.55,2.0) {TypeEncoder};
\node[enc] (eh)  at (2.55,1.0) {HopEncoder};
\node[enc] (eti) at (2.55,0.0) {TimeEncoder};
\node[enc] (ef)  at (2.55,-1.0) {TfsEncoder};
\node[enc] (ep)  at (2.55,-2.0) {GNNPEEncoder};

\node[mod,minimum width=2.0cm]  (fusion) at (5.2,0) {CrossModal\\GatedFusion};
\node[mod,minimum width=2.15cm] (local)  at (7.85,0.88) {LocalModule\\SwiGLU\\DropPath};
\node[mod,minimum width=2.15cm] (global) at (7.85,-0.88) {GlobalModule\\Gumbel\\Codebook};
\node[mod,minimum width=1.85cm] (bridge) at (10.35,0) {CrossAttention\\Bridge};
\node[out,minimum width=1.35cm] (head) at (12.55,0) {MLP\\Head};

\foreach \a/\b in {types/et,hops/eh,times/eti,feat/ef,edge/ep}
  \draw[arr,gray!70] (\a) -- (\b);
\foreach \a in {et,eh,eti,ef,ep}
  \draw[arr,mainblue!70] (\a.east) -- ++(0.28,0) |- (fusion.west);
\draw[arr] (fusion) -- (local);
\draw[arr] (fusion) -- (global);
\draw[arr] (local) -- (bridge);
\draw[arr] (global) -- (bridge);
\draw[arr] (bridge) -- (head);
\end{tikzpicture}
\caption{Luồng module chính trong \RelGTpp{}.}
\end{figure}

\subsection{CrossModalGatedFusion}

Trong \RelGT{} gốc, năm embedding được nối rồi chiếu tuyến tính. \RelGTpp{}
thay bằng cơ chế gate theo phương thức:
\[
  g_i=\sigma(W_\text{self}^{(i)}e_i+
  W_\text{ctx}^{(i)}\text{ctx}_i),
\quad
  \text{ctx}_i=\frac{1}{M-1}\sum_{j\ne i}e_j.
\]
Biểu diễn hợp nhất:
\[
  h_\text{fused}=
  \text{LN}\left(W_\text{out}
  \sum_{i=1}^{M}g_i\odot W_\text{val}^{(i)}e_i\right).
\]

Ý nghĩa: mô hình không buộc feature, type, hop, time và PE có trọng số ngang
nhau. Ví dụ trong tác vụ churn, time có thể quan trọng hơn; trong tác vụ dự
báo sản phẩm, feature tabular và quan hệ có thể nổi bật hơn.

\subsection{SwiGLU Feed-Forward Network}

\RelGTpp{} thay FFN thông thường bằng SwiGLU:
\[
  \text{SwiGLU}(x)=
  W_\text{down}\bigl(\text{SiLU}(W_\text{gate}x)\odot W_\text{up}x\bigr).
\]
SwiGLU thường cho gradient flow tốt hơn FFN GELU đơn giản vì nhánh gate cho
phép mô hình điều tiết từng chiều ẩn theo input.

\subsection{Stochastic Depth}

DropPath bỏ qua một số residual branch trong huấn luyện:
\[
  \text{DropPath}(x;p)=
  \begin{cases}
    x/(1-p), & \text{với xác suất }1-p,\\
    0, & \text{với xác suất }p.
  \end{cases}
\]
Trong mạng sâu, cơ chế này giảm overfitting và giúp ensemble ngầm nhiều độ
sâu khác nhau.

\subsection{Multi-View Readout}

Thay vì chỉ dùng attention pooling, \RelGTpp{} tổng hợp ba góc nhìn:
\begin{align}
  v_\text{attn} &= \sum_j\alpha_jh_j,\\
  v_\text{mean} &= \frac{1}{K-1}\sum_jh_j,\\
  v_\text{max} &= \max_jh_j.
\end{align}
Sau đó:
\[
  h_\text{local}=\text{LN}\bigl(h_i+
  \text{MLP}([v_\text{attn}\concat v_\text{mean}\concat v_\text{max}])\bigr).
\]
Attention nắm quan hệ có chọn lọc, mean giữ bức tranh trung bình, max giữ
tín hiệu nổi bật hiếm.

\subsection{GumbelSoftCodebook}

Codebook phẳng dễ gặp \emph{codebook collapse}: nhiều centroid không được dùng.
\RelGTpp{} thêm Gumbel noise và soft assignment:
\[
  l_{i,b}=-\|q_i-e_b\|_2^2,\quad
  s_{i,b}=\frac{\exp((l_{i,b}+g_{i,b})/\tau)}
  {\sum_{b'}\exp((l_{i,b'}+g_{i,b'})/\tau)}.
\]
Trong đó:
\[
  g_{i,b}=-\log(-\log U_{i,b}),\quad U_{i,b}\sim\text{Uniform}(0,1).
\]
Nhiệt độ được anneal:
\[
  \tau^{(t+1)}=\max(\tau^{(t)}(1-r),\tau_\text{min}).
\]

Để đo mức sử dụng codebook, \RelGTpp{} dùng entropy chuẩn hóa:
\[
  \mathcal{L}_\text{entropy}
  =1-\frac{-\sum_b p_b\log p_b}{\log B}.
\]
Loss bằng 0 khi phân phối centroid đều, tiến tới 1 khi collapse.

\subsection{CrossAttentionBridge}

\RelGT{} gốc nối $h_\text{local}$ và $h_\text{global}$. \RelGTpp{} dùng bridge
hai chiều:
\begin{align}
  \alpha &= \sigma\left(
  \frac{\langle W_q^lh_l,W_k^gh_g\rangle}{\sqrt{d}}\right),&
  \tilde{h}_l &= h_l+\alpha W_v^g h_g,\\
  \beta &= \sigma\left(
  \frac{\langle W_q^gh_g,W_k^lh_l\rangle}{\sqrt{d}}\right),&
  \tilde{h}_g &= h_g+\beta W_v^l h_l.
\end{align}
Sau đó:
\[
  \gamma=\sigma(W_\text{gate}[\tilde{h}_l\concat\tilde{h}_g]),\quad
  h_\text{bridge}=\text{LN}\left(
  W_\text{out}(\gamma\odot\tilde{h}_l+(1-\gamma)\odot\tilde{h}_g)+h_l
  \right).
\]

\subsection{Kỹ thuật triển khai dữ liệu}

Một điểm mạnh của \RelGTpp{} là pipeline dữ liệu theo HDF5:
\begin{itemize}
  \item lưu trước type, index, hop, relative time và edge index của từng đồ
        thị con;
  \item DataLoader đọc theo batch và scatter feature về shape $[B,K,d]$;
  \item encoder bảng chạy theo nhóm node type để giảm số lần forward.
\end{itemize}

Nếu mã hóa từng token riêng lẻ, batch $B=512$, $K=300$ dẫn đến 153.600 lần
encode. Group-by-type giảm xuống khoảng số loại bảng, thường chỉ vài lần
forward lớn.

\subsection{Các vấn đề kỹ thuật cần chú ý}

\begin{warnbox}[Rủi ro triển khai]
\begin{enumerate}
  \item \textbf{Auxiliary loss có thể chưa được cộng vào training loss}: nếu
        entropy/agreement chỉ được ghi log mà không được cộng vào loss, chúng
        không ảnh hưởng đến gradient.
  \item \textbf{Agreement loss dùng \texttt{no\_grad}}: nếu cosine similarity
        được tính trong vùng không gradient, nó không thể kéo local và global
        lại gần nhau.
  \item \textbf{c\_idx trong DDP}: bộ nhớ gán centroid theo node cần đồng bộ
        giữa các GPU, nếu không mỗi rank có thống kê centroid riêng.
  \item \textbf{GNN-PE stochastic}: tăng expressiveness nhưng có thể tạo variance
        trong inference nếu không kiểm soát seed hoặc averaging.
\end{enumerate}
\end{warnbox}

% ============================================================
\section{\RelGTx{}: Hướng Mở rộng Nghiên cứu}
% ============================================================

\subsection{Động lực}

\RelGT{} và \RelGTpp{} đã mạnh, nhưng vẫn còn các hạn chế:
\begin{itemize}
  \item Local attention vẫn all-pair, chưa tận dụng schema để biết cặp token
        nào đáng chú ý.
  \item Time encoder học tương quan thời gian nhưng chưa tách tín hiệu nguyên
        nhân và tương quan giả.
  \item Codebook còn phẳng, chưa phản ánh phân cấp database--table--cluster.
  \item Hop encoder không phân biệt các đường quan hệ khác nhau có cùng độ dài.
  \item Tokenization chưa thích nghi theo tác vụ.
  \item Biểu diễn chưa được khuyến khích bất biến qua các biến đổi schema.
\end{itemize}

\RelGTx{} được đề xuất như một kiến trúc kế tiếp với sáu module.

\subsection{Schema-Guided Sparse Attention}

Với schema $\mathcal{S}=(\mathcal{T},\mathcal{R})$, định nghĩa ma trận reachability:
\[
  M_{ab}^{(r)}=
  \mathbf{1}\bigl[\text{tồn tại đường đi schema độ dài }\le r
  \text{ từ loại }a\text{ đến loại }b\bigr].
\]
Với token $i,j$ có loại $\phi_i,\phi_j$, attention bias:
\[
  b_{ij}=
  \begin{cases}
    0,& M_{\phi_i,\phi_j}^{(r)}=1,\\
    -\infty,& M_{\phi_i,\phi_j}^{(r)}=0.
  \end{cases}
\]
Sparse attention:
\[
  \text{SGSA}(Q,K,V)=
  \text{softmax}\left(\frac{QK^\top}{\sqrt{d_k}}+B+E_\text{rel}\right)V.
\]
Nếu $\rho$ là tỉ lệ cặp token hợp lệ theo schema, chi phí giảm từ
$\mathcal{O}(K^2d)$ xuống $\mathcal{O}(\rho K^2d)$.

\subsection{Causal Temporal Disentanglement}

CTD tách biểu diễn thành thành phần nguyên nhân và tương quan giả:
\[
  h(v)=f_c(z_c(v))+f_s(z_s(v)).
\]
Loss gồm dự báo, bất biến qua môi trường thời gian và độc lập giữa hai phần:
\[
  \mathcal{L}_\text{CTD}=
  \mathcal{L}_\text{pred}
  +\lambda_1\sum_{e\ne e'}
  \left\|\mathbb{E}_{v\in\mathcal{E}_e}z_c(v)
  -\mathbb{E}_{v\in\mathcal{E}_{e'}}z_c(v)\right\|_2^2
  +\lambda_2\|\text{Cov}(z_c,z_s)\|_F^2.
\]
Mục tiêu là làm biểu diễn causal ổn định hơn khi phân phối thời gian thay đổi.

\subsection{Hierarchical Codebook với Schema Priors}

Thay vì codebook phẳng $B$ centroid, HCS dùng ba tầng:
\[
  \text{database}\rightarrow \text{table type}
  \rightarrow \text{intra-table cluster}
  \rightarrow \text{fine centroid}.
\]
Gán tầng 1 là loại bảng:
\[
  c_v^{(1)}=\varphi(v).
\]
Tầng 2:
\[
  c_v^{(2)}=\arg\min_{b\in\mathcal{B}_{c_v^{(1)}}^{(2)}}
  \|\text{proj}^{(2)}(q_v)-e_b^{(2)}\|_2.
\]
Tầng 3:
\[
  c_v^{(3)}=\arg\min_{b\in\mathcal{B}_{c_v^{(2)}}^{(3)}}
  \|\text{proj}^{(3)}(q_v)-e_b^{(3)}\|_2.
\]

Schema prior có thể thêm separation loss:
\[
  \mathcal{L}_\text{sep}=
  \sum_{t\ne t'}\sum_{b\in\mathcal{B}_t}\sum_{b'\in\mathcal{B}_{t'}}
  \max(0,m-\|e_b-e_{b'}\|_2).
\]

\subsection{Relational Positional Encoding}

Hop encoder chỉ biết độ dài đường đi, không biết đường đi đó đi qua quan hệ
nào. RPE mã hóa chuỗi quan hệ:
\[
  \pi(v_i,v_j)=(r_1,r_2,\ldots,r_k).
\]
Mỗi quan hệ có embedding $e_{r_l}$, rồi dùng path encoder:
\[
  h_\pi(v_i,v_j)=
  \text{Transformer}_\text{path}(e_{r_1},\ldots,e_{r_k})_{\text{CLS}}.
\]
Attention được cộng bias:
\[
  \text{Attn-RPE}_{ij}=\text{Attn}_{ij}+W_\text{rpe}h_\pi(v_i,v_j).
\]

\begin{proposition}
RPE biểu đạt mạnh hơn hop encoder vì nó có thể phân biệt hai đường đi cùng
độ dài nhưng khác thứ tự hoặc khác loại quan hệ, ví dụ
$(\text{customer}\rightarrow\text{order}\rightarrow\text{product})$ và
$(\text{product}\rightarrow\text{review}\rightarrow\text{customer})$.
\end{proposition}

\subsection{Meta-Learning Tokenization}

Các tác vụ trong RelBench có nhu cầu khác nhau: task churn cần thời gian,
task sales cần feature sản phẩm, task tương tác cần cấu trúc. MLT sinh bộ
điều biến theo task:
\[
  h_\text{token}^{(\tau)}=
  O\left(g_\tau\odot
  [h_\text{feat}\concat h_\text{type}\concat h_\text{hop}
  \concat h_\text{time}\concat h_\text{pe}]\right),
\]
\[
  g_\tau=\sigma(\text{MLP}_\text{task}(c_\tau)),
\]
trong đó $c_\tau$ mô tả loại tác vụ, số lớp, phân phối nhãn hoặc thống kê
schema liên quan.

\subsection{Contrastive Schema Alignment}

CSA tạo hai view schema-preserving của cùng node, ví dụ column dropout, table
rename, relation permutation hoặc subgraph masking. Loss InfoNCE:
\[
  \mathcal{L}_\text{CSA}=
  -\frac{1}{N}\sum_{i=1}^{N}
  \log
  \frac{\exp(\text{sim}(h_i^{(1)},h_i^{(2)})/\tau_c)}
  {\sum_j\exp(\text{sim}(h_i^{(1)},h_j^{(2)})/\tau_c)}.
\]
Mục tiêu là biểu diễn cùng một thực thể vẫn gần nhau dù schema/view thay đổi.

\subsection{Loss tổng thể của \RelGTx{}}

\[
  \boxed{
  \mathcal{L}_\text{total}=
  \mathcal{L}_\text{task}
  +\lambda_\text{CTD}\mathcal{L}_\text{CTD}
  +\lambda_\text{HCS}\mathcal{L}_\text{sep}
  +\lambda_\text{CSA}\mathcal{L}_\text{CSA}
  +\lambda_\text{entropy}\mathcal{L}_\text{entropy}
  +\lambda_\text{agree}\mathcal{L}_\text{local-global}
  }.
\]

% ============================================================
\section{So sánh Ba Thế hệ}
% ============================================================

\begin{longtable}{p{3.2cm}p{3.7cm}p{3.7cm}p{3.7cm}}
\caption{So sánh \RelGT{}, \RelGTpp{} và \RelGTx{}.}\\
\toprule
\textbf{Khía cạnh} & \textbf{\RelGT{}} & \textbf{\RelGTpp{}} & \textbf{\RelGTx{}} \\
\midrule
\endfirsthead
\toprule
\textbf{Khía cạnh} & \textbf{\RelGT{}} & \textbf{\RelGTpp{}} & \textbf{\RelGTx{}} \\
\midrule
\endhead
Token fusion &
Concat 5 thành phần rồi linear projection. &
CrossModalGatedFusion, gate từng phương thức theo ngữ cảnh. &
Task modulator/meta-learning điều chỉnh trọng số token theo tác vụ. \\
\midrule
Local attention &
All-pair attention trên $K$ token. &
All-pair attention với SwiGLU, DropPath, Multi-View Readout. &
Schema-Guided Sparse Attention, giảm nhiễu và FLOPs. \\
\midrule
Temporal modeling &
Relative time encoder bằng sinusoidal PE. &
Time encoder được gate cùng các modality khác. &
CTD tách causal và spurious temporal signal. \\
\midrule
Positional encoding &
Hop encoder + GNN-PE stochastic. &
GNN-PE cải tiến bằng GIN/residual. &
RPE mã hóa đường đi quan hệ, phân biệt loại quan hệ. \\
\midrule
Global memory &
Flat centroid/codebook. &
GumbelSoftCodebook, entropy tracking, temperature annealing. &
Hierarchical codebook theo schema/table/cluster. \\
\midrule
Local-global fusion &
Concat + MLP. &
CrossAttentionBridge hai chiều. &
Bridge nâng cao kết hợp causal và hierarchical global context. \\
\midrule
Regularization &
Dropout/weight decay cơ bản. &
DropPath, entropy, agreement loss nếu được kích hoạt. &
CSA, CTD, HCS separation, entropy schedule. \\
\midrule
Transfer giữa task &
Không phải trọng tâm. &
Không phải trọng tâm. &
Meta-learning tokenization và schema alignment. \\
\midrule
Độ phức tạp local &
$\mathcal{O}(K^2d)$. &
$\mathcal{O}(K^2d)$, nhưng module ổn định hơn. &
$\mathcal{O}(\rho K^2d)$ với $\rho<1$. \\
\bottomrule
\end{longtable}

\subsection{Bảng điểm mạnh và điểm yếu}

\begin{table}[H]
\centering\small
\caption{Đánh giá tổng hợp.}
\begin{tabularx}{\textwidth}{lXX}
\toprule
\textbf{Phiên bản} & \textbf{Điểm mạnh} & \textbf{Điểm yếu / rủi ro} \\
\midrule
\RelGT{} &
Thiết kế rõ ràng, token hóa giàu ngữ cảnh, kết hợp local-global, phù hợp
dữ liệu quan hệ có thời gian. &
Fusion còn tuyến tính; attention chưa biết schema; codebook phẳng; GNN-PE
có variance. \\
\RelGTpp{} &
Mạnh hơn về kỹ thuật huấn luyện; gate đa phương thức; SwiGLU; codebook mềm;
readout nhiều góc nhìn. &
Auxiliary loss cần kiểm tra; DDP/codebook sync có thể phức tạp; số tham số
và pipeline tăng. \\
\RelGTx{} &
Đưa prior schema, causal temporal, hierarchical memory và transfer learning
vào kiến trúc. &
Nhiều module mới làm tăng độ khó ablation; CTD/CSA cần thiết kế dữ liệu và
schedule cẩn thận. \\
\bottomrule
\end{tabularx}
\end{table}

% ============================================================
\section{Kế hoạch Thực nghiệm Đề xuất}
% ============================================================

\subsection{Ablation theo tầng kiến trúc}

\begin{table}[H]
\centering
\caption{Ablation cho ba thế hệ.}
\begin{tabularx}{\textwidth}{clX}
\toprule
\textbf{\#} & \textbf{Cấu hình} & \textbf{Câu hỏi kiểm tra} \\
\midrule
1 & GNN/LightGBM baseline & RELGT có thật sự vượt mô hình tabular/GNN không? \\
2 & \RelGT{} full & Token hóa 5 phần + local-global giúp bao nhiêu? \\
3 & \RelGT{} bỏ global & Global centroid đóng góp bao nhiêu? \\
4 & \RelGT{} bỏ GNN-PE/time/hop & Thành phần token nào quan trọng nhất theo task? \\
5 & \RelGTpp{} full & Các cải tiến engineering có tăng ổn định/hiệu suất không? \\
6 & \RelGTpp{} bỏ Gumbel/Bridge/Fusion & Từng module mới đóng góp bao nhiêu? \\
7 & \RelGTx{} thêm SGSA & Schema mask giảm nhiễu và FLOPs thế nào? \\
8 & \RelGTx{} thêm RPE/CTD/HCS/MLT/CSA & Kiểm tra lần lượt biểu đạt quan hệ, nhân quả, memory và transfer. \\
\bottomrule
\end{tabularx}
\end{table}

\subsection{Chỉ số cần báo cáo}

Ngoài metric chính của task, nên báo cáo thêm:
\begin{itemize}
  \item \textbf{Codebook utilization}: entropy, số centroid được dùng, histogram
        gán centroid.
  \item \textbf{Attention sparsity}: tỉ lệ $\rho$, FLOPs thực tế, memory peak.
  \item \textbf{Temporal robustness}: hiệu suất theo split thời gian, đặc biệt
        khi test nằm xa train.
  \item \textbf{Variance}: mean/std qua nhiều seed, nhất là với GNN-PE stochastic.
  \item \textbf{Ablation theo schema}: hiệu suất khi đổi hoặc mask một phần schema.
\end{itemize}

\subsection{Cấu hình tham khảo}

\begin{table}[H]
\centering
\caption{Siêu tham số tham khảo.}
\begin{tabular}{lcc}
\toprule
\textbf{Tham số} & \textbf{Giá trị smoke/test} & \textbf{Giá trị đầy đủ} \\
\midrule
$K$ token/seed & 64--128 & 300 \\
Embedding dim $d$ & 64--128 & 512 \\
Centroid $B$ & 256--512 & 4096 \\
Attention heads & 4 & 4--8 \\
Transformer layers & 1 & 1--3 \\
Learning rate & $10^{-4}$ & $10^{-4}$ \\
Weight decay & $10^{-5}$ & $10^{-5}$ \\
Dropout & 0.1 & 0.1 \\
Gradient clipping & 1.0 & 1.0 \\
\bottomrule
\end{tabular}
\end{table}

% ============================================================
\section{Thảo luận}
% ============================================================

\subsection{Ý nghĩa khoa học}

\RelGT{} cho thấy Graph Transformer có thể dùng hiệu quả trên dữ liệu quan hệ
nếu bài toán tokenization được thiết kế đúng. Đóng góp không nằm ở việc áp
dụng Transformer một cách trực tiếp, mà ở cách biến node trong cơ sở dữ liệu
quan hệ thành token có đủ năm lớp thông tin: nội dung, kiểu, khoảng cách,
thời gian và vị trí.

\RelGTpp{} nhấn mạnh một điểm khác: với kiến trúc phức tạp, hiệu quả thực tế
phụ thuộc nhiều vào các chi tiết kỹ thuật như fusion, readout, codebook update,
loss phụ và pipeline dữ liệu. Một ý tưởng tốt có thể yếu nếu triển khai không
ổn định; ngược lại, một loạt cải tiến nhỏ nhưng đúng chỗ có thể làm mô hình
dễ train hơn đáng kể.

\RelGTx{} mở rộng từ ``học trên graph quan hệ'' sang ``học có prior về schema
và thời gian''. Đây là bước tự nhiên vì cơ sở dữ liệu quan hệ không chỉ là
graph, mà là graph được sinh ra bởi schema có ý nghĩa.

\subsection{Giới hạn}

Các hướng mở rộng trong \RelGTx{} cần được kiểm chứng cẩn thận. SGSA phụ thuộc
vào chất lượng schema mask; CTD cần định nghĩa môi trường thời gian hợp lý;
HCS cần tránh phân mảnh codebook; CSA cần augmentation không phá ngữ nghĩa.
Do đó, \RelGTx{} nên được phát triển theo từng module và luôn có ablation.

\subsection{Kết luận}

\RelGT{} là một bước tiến quan trọng cho Relational Deep Learning vì nó đưa
Transformer vào dữ liệu quan hệ bằng một thiết kế token hóa phù hợp. \RelGTpp{}
củng cố kiến trúc đó bằng các cải tiến thực dụng về fusion, readout, codebook
và huấn luyện. \RelGTx{} là định hướng nghiên cứu tiếp theo, khai thác sâu
hơn schema, nhân quả thời gian và khả năng transfer giữa tác vụ.

Thông điệp chính của toàn bộ chuỗi này là: để học tốt trên cơ sở dữ liệu quan
hệ, mô hình cần hiểu đồng thời \textbf{bảng nào}, \textbf{quan hệ nào},
\textbf{khi nào}, \textbf{ở đâu trong đồ thị con}, và \textbf{ngữ cảnh toàn
cục nào} đang ảnh hưởng đến dự báo.

% ============================================================
\begin{thebibliography}{99}

\bibitem{relgt2025}
Dwivedi, V. P., Jaladi, S., Shen, Y., López, F., Kanatsoulis, C. I.,
Puri, R., Fey, M., \& Leskovec, J. (2025).
\textit{Relational Graph Transformer}.
arXiv:2505.10960.

\bibitem{relbench2023}
Robinson, J. et al. (2023).
\textit{RelBench: A Benchmark for Deep Learning on Relational Databases}.
NeurIPS.

\bibitem{attention2017}
Vaswani, A. et al. (2017).
\textit{Attention Is All You Need}.
NeurIPS.

\bibitem{graphsage2017}
Hamilton, W., Ying, Z., \& Leskovec, J. (2017).
\textit{Inductive Representation Learning on Large Graphs}.
NeurIPS.

\bibitem{hgt2020}
Hu, Z. et al. (2020).
\textit{Heterogeneous Graph Transformer}.
WWW.

\bibitem{graphgps2022}
Rampášek, L. et al. (2022).
\textit{Recipe for a General, Powerful, Scalable Graph Transformer}.
NeurIPS.

\bibitem{gin2019}
Xu, K. et al. (2019).
\textit{How Powerful are Graph Neural Networks?}
ICLR.

\bibitem{gumbel2017}
Jang, E., Gu, S., \& Poole, B. (2017).
\textit{Categorical Reparameterization with Gumbel-Softmax}.
ICLR.

\bibitem{maml2017}
Finn, C., Abbeel, P., \& Levine, S. (2017).
\textit{Model-Agnostic Meta-Learning for Fast Adaptation of Deep Networks}.
ICML.

\bibitem{simclr2020}
Chen, T. et al. (2020).
\textit{A Simple Framework for Contrastive Learning of Visual Representations}.
ICML.

\bibitem{causalrep2021}
Schölkopf, B. et al. (2021).
\textit{Toward Causal Representation Learning}.
Proceedings of the IEEE.

\end{thebibliography}

\end{document}
